\chapter{أدوات البناء}
\label{ch:build}

تتعرف على الأدوات التي تنقل بها عمل مؤسستك الخاص إلى كيان: ماذا تبني بكل أداة، وبأي ترتيب، ولماذا لا يصل أي تعديل إلى الفريق قبل نشره.

\section{لماذا أدوات البناء؟}
\label{sec:build-why}
\begin{rcslide}{لماذا أدوات البناء؟}
\sphase{الفكرة}
\slead{لكل مؤسسة أعمال لا تشبه غيرها. وبدل أن تنتظر تطويراً خاصاً، يبنيها مسؤول مخوَّل بالاختيار والتعبئة.}
\scompare{في برنامج تقليدي}{تطلب ميزة من فريق البرمجة \\ تنتظر أسابيع أو أشهراً \\ كل تعديل طلب جديد}{في كيان}{يبني المسؤول المخوَّل ما يحتاجه \\ يجرّبه وينشره بنفسه \\ يعدّله كلما تغيّر العمل}
\stake{لا تحتاج أدوات البناء إلى كتابة شيفرة، لكنها تحتاج إلى فهم جيد للعمل نفسه. لذلك يشارك فيها من يعرف العمل.}
\slink{builder}{دليل أدوات البناء}{ar/builder.html}
\end{rcslide}

\section{الكيان: نوع السجلات}
\label{sec:build-entity}
\begin{rcslide}{الكيان: نوع السجلات}
\sphase{الأساس}
\slead{الكيان هو أساس كل ما تبنيه: يحدد المعلومات التي تُحفظ، وطريقة عرضها، ومن يعمل عليها.}
\squote{ما الذي يمثله السجل الواحد؟ منشأة؟ مستفيد؟ طلب؟ عقد؟}
\scard{form}{الحقول}{المعلومات التي يحفظها كل سجل.}
\scard{paper}{الصفحة والقوائم}{كيف يظهر السجل، وكيف تُعرض السجلات معاً.}
\scard{flow}{القواعد}{الحالات وما يُحسب تلقائياً.}
\scard{lock}{الصلاحيات}{من يرى السجل، ومن يعدّله.}
\stake{كل أداة أخرى تعمل على كيان: النموذج يكتب فيه، والإجراء يبدأ منه، واللوحة تحسب منه. لذلك تبدأ به.}
\slink{builder}{الكيانات}{ar/builder-entities.html}
\end{rcslide}

\section{مسودة ثم نشر}
\label{sec:build-publish}
\begin{rcslide}{مسودة ثم نشر}
\sphase{الأساس}
\slead{كل تعديل تجريه يبقى مسودة لا يراها غيرك، حتى تراجعه وتنشره.}
\schain{تبني ← تجرّب بسجل فعلي ← تراجع التغييرات ← تنشر ← يراه الفريق}
\scard{lock}{تجربة آمنة}{تخطئ وتصحح دون أن يتأثر عمل أحد.}
\scard{people}{الفريق}{يعمل دائماً على الإصدار المنشور.}
\stake{القاعدة نفسها في الكيانات والنماذج والإجراءات: لا يصل شيء إلى الفريق قبل النشر.}
\slink{builder}{المراجعة والنشر}{ar/builder-publish.html}
\end{rcslide}

\section{التطبيقات والتنقل}
\label{sec:build-nav}
\begin{rcslide}{التطبيقات والتنقل}
\sphase{الأساس}
\slead{لا يكفي أن تنشر الكيان؛ يجب أن يجده الأعضاء في المكان الذي يتوقعونه.}
\schain{التطبيق ← القسم ← الصفحة}
\scard{building}{التطبيقات والتنقل}{قائمة التطبيقات وأقسامها وهوية كل تطبيق.}
\scard{tasks}{التنقل حسب العمل}{ما يراه كل عضو أولاً، مثل قائمة «المستحق عليّ».}
\stake{يجد كل موظف عمله في مكان واضح، بدل أن يبحث عنه بين صفحات لا تخصه.}
\slink{builder}{التطبيقات والتنقل}{ar/builder-navigation.html}
\slink{builder}{التنقل حسب العمل}{ar/builder-work-navigation.html}
\end{rcslide}

\section{النماذج}
\label{sec:build-forms}
\begin{rcslide}{النماذج}
\sphase{جمع البيانات}
\slead{النموذج صفحة أسئلة يعبّئها الموظف أو مقدّم الطلب، وتصل إجاباته إلى من يراجعها، وقد تُحفظ في سجل.}
\scompare{الاستمارة الورقية}{تُعبّأ بخط اليد \\ تصل ناقصة أحياناً \\ تُنقل يدوياً إلى جدول}{النموذج في كيان}{يُعبّأ من الحاسوب أو الجوال \\ قواعد تمنع الإجابة الناقصة \\ تُحفظ الإجابات في السجل مباشرة}
\scard{flow}{مسار الطلب}{من يراجع الطلب، وبأي ترتيب يوافق.}
\scard{lock}{الوصول إلى الردود}{من يرى الإجابات، ومن يعمل عليها.}
\stake{ابدأ من الاستمارات الورقية التي تستخدمها مؤسستك اليوم؛ فهي أسرع ما تنقله إلى كيان.}
\slink{builder}{النماذج}{ar/builder-forms.html}
\slink{inspection}{تطبيق عملي: الكشف على المنشآت}{ar/index.html}
\end{rcslide}

\section{الخدمات}
\label{sec:build-services}
\begin{rcslide}{الخدمات}
\sphase{جمع البيانات}
\slead{الخدمة تضع نموذجاً أو أكثر أمام المتعاملين على بوابة المؤسسة، باسم واضح ودليل يقرؤونه قبل البدء.}
\scard{paper}{التعريف والدليل}{ما الخدمة، وما المطلوب قبل التقديم.}
\scard{people}{الأهلية}{من يحق له التقديم.}
\scard{sheet}{الرسوم}{ما يدفعه مقدّم الطلب.}
\scard{flow}{المسار والنماذج}{النماذج التي يعبّئها، والطريق الذي يسلكه طلبه.}
\stake{ابنِ النموذج وانشره أولاً، ثم أنشئ الخدمة حوله.}
\slink{builder}{الخدمات}{ar/builder-services.html}
\end{rcslide}

\section{الإجراءات الداخلية}
\label{sec:build-isop}
\begin{rcslide}{الإجراءات الداخلية}
\sphase{أتمتة العمل}
\slead{مسار من خطوات يبدأ عند حدث تحدده، فلا يعتمد إنجاز العمل على أن يتذكر أحد الخطوة التالية.}
\schain{حدث: إرسال نموذج ← موافقة ← مهمة ← إشعار}
\scard{flow}{المشغّل}{الحدث الذي يبدأ الإجراء، مثل إنشاء سجل.}
\scard{tasks}{الخطوات}{موافقة، أو مهمة، أو إشعار.}
\scard{chart}{التشغيل}{تتابع كل تشغيل وما حدث فيه.}
\stake{تبنيه، ثم تختبره، ثم تنشره، ثم تتابع تشغيله. هذه دورة كل إجراء.}
\slink{builder}{الإجراءات الداخلية}{ar/builder-internal-sops.html}
\slink{sop}{تطبيق عملي: الإجراءات الداخلية}{ar/index.html}
\end{rcslide}

\section{الأتمتة والبوتات}
\label{sec:build-automation}
\begin{rcslide}{الأتمتة والبوتات}
\sphase{أتمتة العمل}
\slead{أداتان تربطان كيان بما حوله: مسارات تعمل وحدها، ومساعدون يعملون مع الأعضاء.}
\scard{flow}{إجراءات التشغيل القياسية}{مسار يبدأ بحدث في كيان، ثم يقرأ البيانات ويتحقق منها وينفّذ إجراءً.}
\scard{link}{الاتصالات}{ربط كيان بخدمات وحسابات أخرى.}
\scard{chat}{البوتات}{مساعدون باسم وصورة، يضيفهم الأعضاء في الدردشة والبريد والتقويم.}
\stake{كل إجراء ينفّذه المسار يترك أثراً تستطيع مراجعته.}
\slink{builder}{إجراءات التشغيل القياسية}{ar/builder-sops.html}
\slink{builder}{البوتات}{ar/builder-bots.html}
\end{rcslide}

\section{المستندات والطباعة}
\label{sec:build-docs}
\begin{rcslide}{المستندات والطباعة}
\sphase{المخرجات}
\slead{مستند رسمي يُملأ من بيانات السجل، فلا تعيد كتابة خطاب أو كشف في كل مرة.}
\scard{paper}{قوالب Word}{تكتب القالب في Word وترفعه.}
\scard{sheet}{قوالب Excel}{للكشوف والجداول.}
\scard{form}{قوالب الطباعة}{تصممها داخل كيان دون ملف خارجي.}
\scard{building}{الهويات والترويسات}{شعار المؤسسة وبياناتها ورأس الصفحة وتذييلها.}
\stake{إذا كانت مستنداتك رسمية، فابدأ بالهوية والترويسة، ثم اختر نوع القالب.}
\slink{builder}{قوالب DOCX}{ar/builder-docx-templates.html}
\slink{builder}{محرر الترويسة}{ar/builder-letterheads.html}
\end{rcslide}

\section{لوحات المعلومات والتقارير}
\label{sec:build-dash}
\begin{rcslide}{لوحات المعلومات والتقارير}
\sphase{المخرجات}
\slead{بعد أن تتراكم السجلات، تلخّص اللوحة العمل بأرقام: كم طلباً مفتوحاً؟ كم متأخراً؟}
\schain{تنشئ اللوحة ← تضيف العناصر ← تربطها بالبيانات ← تعاينها ← تنشرها}
\scard{chart}{العناصر والمرشحات}{أرقام ومخططات يصفّيها المشاهد بنفسه.}
\scard{ai}{بالذكاء الاصطناعي}{تصف اللوحة أو تسأل سؤالك، فيُبنى الجواب من بيانات المؤسسة.}
\scard{lock}{بحسب الصلاحيات}{يرى كل مشاهد ما تسمح به صلاحياته.}
\stake{اللوحة لا تحتاج إلى تجميع يدوي؛ أرقامها من السجلات نفسها.}
\slink{builder}{لوحات المعلومات}{ar/builder-dashboards.html}
\slink{builder}{تقارير الذكاء الاصطناعي}{ar/builder-ai-reports.html}
\end{rcslide}

\section{ترتيب البناء}
\label{sec:build-order}
\begin{rcslide}{ترتيب البناء}
\sphase{الخلاصة}
\slead{كل أداة تعتمد على التي قبلها، فابنِ بهذا الترتيب وأضف ما يحتاجه العمل فقط.}
\schain{الكيان ← النشر والتنقل ← النموذج ← الإجراء ← المستند ← اللوحة}
\smetric{أول ما تبنيه}{نموذج بلا سجل}{كيان تُحفظ فيه البيانات}
\smetric{قبل أن يراه الفريق}{تعديل مباشر}{مسودة ثم نشر}
\smetric{متى تضيف أداة؟}{كل الأدوات دفعة واحدة}{حين يحتاجها العمل}
\stake{رأيت الأدوات منفصلة. وفي التطبيقات العملية تراها تعمل معاً في مسار كامل.}
\slink{inspection}{تطبيق عملي: الكشف على المنشآت}{ar/index.html}
\slink{sop}{تطبيق عملي: الإجراءات الداخلية}{ar/index.html}
\end{rcslide}
